% Part: turing-machines
% Chapter: machines-computations
% Section: well-behaved-machines

\documentclass[../../../include/open-logic-section]{subfiles}

\begin{document}

\olfileid{tur}{mac}{dis}
\olsection{શિસ્તબદ્ધ મશીનો}

\begin{explain}
\olref[hal]{sec} વિભાગમાં આપણે એવી ટ્યુરિંગ મશીનો
જોઈ હતી જેમાં એક જ નિર્ધારિત થંભવાની અવસ્થા~$h$
હોય છે---આવાં મશીનો થંભે તો અવસ્થા~$h$માં જ થંભે
તેની ખાતરી હોય છે. આ રીતે, એક જ થંભવાની અવસ્થાવાળાં
મશીનો સામાન્ય ટ્યુરિંગ મશીનોને આપણે આપેલી છૂટ કરતાં
વધુ ``શિસ્તબદ્ધ'' છે. ટ્યુરિંગ મશીનોના વર્તન પર બીજી
શરતો પણ મૂકી શકીએ. દાખલા તરીકે, મશીન ખાનું~$0$માંનું
ટેપના છેડાનું ચિહ્ન ક્યારેય ભૂંસે નહીં અથવા ખાનું~$0$માંથી
ડાબે જવાનો પ્રયત્ન ન કરે એવી મનાઈ આપણે મૂકી નથી.
(આવા કિસ્સામાં હેડ ખાનું~$0$માં જ રહે છે એમ આપણી
વ્યાખ્યા કહે છે; બીજી વ્યાખ્યાઓમાં મશીન થંભે છે.)
એવી જ રીતે, ટ્યુરિંગ મશીન થંભે તો ખાનું~$1$ વાંચતાં
થંભે એવી ધારણા કરી શકાય તે ક્યારેક ઇચ્છનીય છે.
\end{explain}

\begin{defn}\ollabel{defn:disciplined}
ટ્યુરિંગ મશીન~$M$ \emph{શિસ્તબદ્ધ} છે ત્યારે અને
ત્યારે જ જ્યારે
\begin{enumerate}
    \item તેમાં એક જ નિર્ધારિત થંભવાની અવસ્થા~$h$ છે,
    \item તે થંભે તો ખાનું~$1$ વાંચતાં થંભે છે,
    \item તે ખાનું~$0$માંનું $\TMendtape$ પ્રતીક ક્યારેય
    ભૂંસતું નથી, અને
    \item તે ખાનું~$0$માંથી ડાબે જવાનો પ્રયત્ન ક્યારેય
    કરતું નથી.
\end{enumerate}
\end{defn}

\begin{explain}
આપણે અગાઉ ચર્ચ્યું છે કે કોઈ પણ ટ્યુરિંગ મશીનને એ જ
વર્તન ધરાવતું પરંતુ નિર્ધારિત થંભવાની અવસ્થાવાળું મશીન
બનાવી શકાય છે. આ માટે માત્ર નવી અવસ્થા~$h$ ઉમેરીને
મૂળ $\delta$ જે જોડી $\tuple{q,\sigma}$ પર અવ્યાખ્યાયિત
હોય તે દરેક માટે સૂચના
$\delta(q, \sigma) = \tuple{h, \sigma, N}$ ઉમેરવી પડે.
સાબિતી કંટાળાજનક હોવા છતાં એ સાચું છે કે કોઈ પણ
ટ્યુરિંગ મશીન~$M$ને એવા શિસ્તબદ્ધ ટ્યુરિંગ મશીન~$M'$
માં ફેરવી શકાય છે જે એ જ ઇનપુટ પર થંભે છે અને એ જ
આઉટપુટ આપે છે. દાખલા તરીકે, ટ્યુરિંગ મશીન થંભે અને
તે ખાનું~$1$ પર ન હોય, તો હેડ ટેપના છેડાનું ચિહ્ન
મળે ત્યાં સુધી ડાબે ખસે, પછી એક ખાનું જમણે ખસે અને
પછી થંભે એવી સૂચનાઓ ઉમેરી શકાય. બીજી શરતો કેવી
રીતે પૂરી કરી શકાય તે તમે વિચારો.
\end{explain}

\begin{ex}
    \olref{fig:adder-disc}માં આપણે
    \olref[una]{ex:adder}ના સરવાળાના મશીનને શિસ્તબદ્ધ
    મશીનમાં ફેરવીએ છીએ.
    \begin{figure}\[
\begin{tikzpicture}[->,>=stealth',shorten >=1pt,auto,node distance=2.8cm,
                    semithick]
  \tikzstyle{every state}=[fill=none,draw=black,text=black]

  \node[initial,state]         (A)              {$q_0$};
  \node[state]         (B) [right of=A] {$q_1$};
  \node[state]         (C) [below right of=B] {$q_2$};
  \node[state]         (D) [below left of=C] {$q_3$};
  \node[state]         (H) [left of=D] {$h$};

  \path (A) edge node {\TMtrans{\TMblank}{\TMstroke}{\TMstay}} (B)
           edge [loop below] node {\TMtrans{\TMstroke}{\TMstroke}{\TMright}} (B)
        (B) edge [loop below] node {\TMtrans{\TMstroke}{\TMstroke}{\TMright}} (B)
            edge node {\TMtrans{\TMblank}{\TMblank}{\TMleft}} (C)
        (C) edge node {\TMtrans{\TMstroke}{\TMblank}{\TMleft}} (D)
        (D) edge [loop above] node {\TMtrans{\TMstroke}{\TMstroke}{\TMleft}} (D)
        edge node {\TMtrans{\TMendtape}{\TMendtape}{\TMright}} (H);
\end{tikzpicture}
\]\caption{શિસ્તબદ્ધ સરવાળાનું મશીન}
\ollabel{fig:adder-disc}
\end{figure}
\end{ex}

\begin{prop}\ollabel{prop:disciplined} દરેક ટ્યુરિંગ
મશીન~$M$ માટે એવું શિસ્તબદ્ધ ટ્યુરિંગ મશીન~$M'$
છે કે $M$ આઉટપુટ~$O$ સાથે થંભે તો તે પણ
આઉટપુટ~$O$ સાથે થંભે, અને $M$ ન થંભે તો તે
પણ ન થંભે. ખાસ કરીને, ટ્યુરિંગ મશીન વડે સંગણનીય
દરેક વિધેય $f\colon\Nat^n \to \Nat$ શિસ્તબદ્ધ
ટ્યુરિંગ મશીન વડે પણ સંગણનીય છે.
\end{prop}

\begin{prob}\label{tur:mac:dis:prob:disc-succ}
$f(x) = x+1$ ગણતું શિસ્તબદ્ધ મશીન આપો.
\end{prob}


\begin{prob}\label{tur:mac:dis:prob:copier}
ઇનપુટ $\TMstroke^n$ પર શરૂ કરતાં આઉટપુટ
$\TMstroke^n \concat \TMblank \concat \TMstroke^n$
આપે એવું શિસ્તબદ્ધ મશીન શોધો.
\end{prob}

\end{document}
